% Elaborado por Fabian Gomez
% Version 1.0

\section{References}

Esta sección lista los documentos normativos y guías técnicas que enmarcan la estructura, el contenido mínimo y las buenas prácticas aplicadas en este SyRS/SRS del Rover Olympus \cite{ieee_29148, segoldmine_ppi}. Los estándares principales son ISO/IEC/IEEE 29148:2018 (proceso de ingeniería de requerimientos), el NASA Systems Engineering Handbook (prácticas SE y tailoring) y el CubeSat Space Protocol (encapsulamiento de telemetría/comandos en la interfaz rover–estación terrena) \cite{ieee_29148, nasa_se_handbook, csp_spec}. Se complementan con estándares ECSS y NASA aplicables al análisis de fallas, calidad de software y protección planetaria conforme se detalla en \S\ref{subsec:applicable_standards} \cite{ecss_e_st_10_06c, ecss_q_st_30, nasa_std_8739_8, nasa_planetary_protection}.

\subsection{Applicable Standards}\label{subsec:applicable_standards}

Los siguientes estándares y guías se consideran aplicables o de referencia directa para la elaboración, auditoría y mantenimiento del SyRS/SRS. Se distingue entre documentos \emph{principales} (estructura y metodología del documento) y \emph{de apoyo} (análisis específicos: FMEA, calidad de software, protección planetaria).

\subsubsection*{Documentos principales}

\begin{enumerate}
    \item \textbf{ISO/IEC/IEEE 29148:2018 – Systems and software engineering — Life cycle processes — Requirements engineering.} \cite{ieee_29148}
    \begin{itemize}
        \item Aplica como marco para estructura del documento, criterios de ``buen requisito'', contenido esperado y prácticas de gestión/ingeniería de requisitos.
    \end{itemize}

    \item \textbf{NASA Systems Engineering Handbook – NASA SP‑2016‑6105 Rev2.} \cite{nasa_se_handbook}
    \begin{itemize}
        \item Aplica como guía de buenas prácticas SE (incluyendo relación requisitos\allowbreak –\allowbreak arquitectura\allowbreak –\allowbreak V\&V, y recomendaciones de tailoring para el contexto de proyecto).
    \end{itemize}

    \item \textbf{CubeSat Space Protocol (CSP) Specification.} \cite{csp_spec}
    \begin{itemize}
        \item Aplica como referencia técnica para el encapsulamiento y verificación de integridad de paquetes de telemetría/comandos en la interfaz rover–estación terrena (ver \S\ref{sec:system_interfaces}).
    \end{itemize}
\end{enumerate}

\subsubsection*{Documentos de apoyo}

\begin{enumerate}
    \item \textbf{ECSS-E-ST-10-06C – Technical Requirements Specification.} \cite{ecss_e_st_10_06c}
    \begin{itemize}
        \item Aplica como referencia para el formato y atributos de requerimientos técnicos en proyectos espaciales europeos; consultada en la auditoría 2026-04-04 para validar completitud y verificabilidad.
    \end{itemize}

    \item \textbf{ECSS-Q-ST-30C – Dependability.} \cite{ecss_q_st_30}
    \begin{itemize}
        \item Aplica como marco del análisis FMEA (Anexo, \S\ref{sec:appendices}): clasificación severidad/ocurrencia/detección requerida para la cobertura RF-005 (Fault Stop) y RNF-002 (homing).
    \end{itemize}

    \item \textbf{NASA-STD-8739.8B – Software Assurance and Software Safety.} \cite{nasa_std_8739_8}
    \begin{itemize}
        \item Aplica como referencia para la clasificación de criticidad del software de control del rover y los criterios de aseguramiento.
    \end{itemize}

    \item \textbf{NASA Planetary Protection Provisions for Robotic Extraterrestrial Missions.} \cite{nasa_planetary_protection}
    \begin{itemize}
        \item Aplica como referencia consultada para el ajuste del UC-05 (\textit{Safe Deactivation / Passivation}) en el contexto de no interferencia con futuras misiones lunares de exploración.
    \end{itemize}
\end{enumerate}
